% !PRINTED 628
\weberpage{628}
%<<CONT
% [sic] The source defines the resolvent as F(v)=0 but subsequently prints F(z) here; notation retained source-faithfully.
common divisor than the identical permutation. For every common divisor of all these groups would have to be a normal divisor of the symmetric, and consequently also of the alternating group; but since it cannot be the alternating group itself, it must by \S.\,177 reduce to the unit group.

The resolvent $F(v) = 0$ is therefore, according to the designation introduced in \S.\,156, a \gesperrt{total resolvent} of $f(x) = 0$, and its group must be of the same degree as the group of $f(x) = 0$, whose degree is $1 . 2 . 3 . 4 . 5 = 120$. Since $F(z)$ is also irreducible, the permutation group among the indices of the six quantities $v$ is transitive. It is of degree 120, while the degree of the symmetric group of the permutations of six figures is $6 . 120$. With this we have the theorem:

\gesperrt{The symmetric permutation group of six figures has a transitive divisor of index 6.}

In order to find this remarkable group, which we shall denote by $C$, we have only to investigate the influence which all the 120 permutations of the $x$ exert on the $v$, or on the quantities (11) of the preceding section.

Now we have seen earlier (\S.\,153, 2.) that the transpositions $(0, 1)$, $(0, 2)$, $(0, 3)$, $(0, 4)$ generate, by repeated composition, the whole symmetric group of the permutations of five figures. It therefore suffices to determine the influence of these four transpositions on the indices of the $v$. By the application of a transposition $(0, 1)$ each of the expressions (11), \S.\,181 passes into one of the expressions (12) and conversely, and if one determines this change for the four transpositions named, one obtains the following four generating permutations $\pi_1, \pi_2, \pi_3, \pi_4$ of the group $C$

\[
\begin{aligned}
(0, 1)&: \quad \pi_1 = (1, 3) (2, 5) (4, 6) \\
(0, 2)&: \quad \pi_2 = (1, 4) (2, 3) (5, 6) \\
(0, 3)&: \quad \pi_3 = (1, 5) (2, 6) (3, 4) \\
(0, 4)&: \quad \pi_4 = (1, 6) (2, 4) (3, 5),
\end{aligned}
\]

where the transpositions $(1, 3) \ldots$ occurring in the $\pi$ refer to the indices of the $u$ or of the $v$.

Among the 120 permutations of the group $C$ there occur, among others, also
%CONT>>
